\documentclass[11pt]{article}
\usepackage[margin=1in]{geometry}
\usepackage[utf8]{inputenc}
\usepackage{hyperref}
\usepackage{amsmath,amssymb}
\usepackage{fancyhdr}
\usepackage{enumitem}
\pagestyle{fancy}
\fancyhf{}
\fancyfoot[C]{\thepage}
\title{Homotopy transfers of curved \$L\_\textbackslash{}infty\$ algebras}
\author{Shuhan Jiang}
\date{}
\begin{document}
\maketitle

\noindent\textbf{Source:} \href{https://arxiv.org/abs/2606.23965}{arXiv:2606.23965} \quad \textbf{Field:} math.DG\\

\section*{TL;DR}
The paper establishes that L\ensuremath{\infty} spaces over a dg manifold form a category of fibrant objects (CFO), with the stated corollary that transitive L\ensuremath{\infty} algebroids inherit this structure via the companion paper equivalence. The technical\_correctness specialist (confidence 0.70, overall: mostly\_sound) found the supporting scaffolding --- Lemmas 24, 43; Propositions 23, 44; Lemma 17; Lemma 41 --- fully supported on detailed inspection. The headline Theorem 1/30 was assessed as partially\_supported at major severity not because the mathematical reasoning is defective but because no machine-checkable proof artifact accompanies it. The reproducibility specialist (confidence 0.87, score 0.24) independently flags the same gap at critical severity. These two specialist findings trigger the recommendation gate: the field is math.* (code-amenable), and both technical\_correctness and reproducibility flag a missing proof-as-code for the headline claim at major or critical severity, so the default recommendation is major\_revision. Two additional cross-cutting concerns compound this: (a) Theorem 1's abstract corollary and the intermediate Lemma 16 both depend on the companion paper cattaneojiang26, which carries a 2026 date and no arXiv identifier and was therefore unavailable for any reviewer; (b) the novelty specialist (confidence 0.62, score 0.45, verdict: incremental) identified three bodies of closely related prior art absent from the bibliography. The citation specialist timed out entirely before producing output; this is a pipeline failure and not evidence of a paper-level bibliography deficiency. Specialists are not in disagreement on substantive points --- technical\_correctness and reproducibility align on the proof-as-code gap, and novelty's incremental verdict is consistent with technical\_correctness finding a sound but not transformative extension of existing frameworks.

\textbf{Recommendation:} Major revision \quad \textbf{Confidence:} 72\%\\

\section*{Strengths}
\begin{itemize}[leftmargin=*]
  \item The geometric-to-algebraic reduction via the global sections functor \ensuremath{\Gamma} is clean and rigorously supported: Proposition 18's proof correctly chains Proposition 15 and Lemma 13 to establish full faithfulness and the quasi-inverse equivalence, with the finite-rank projectivity condition doing essential work throughout.
  \item The homotopy transfer theorem for curved L\ensuremath{\infty} algebras over filtered cdgas (Proposition 44) is developed with explicit algebraic detail --- the Berglund symmetrized tensor homotopy identity (Lemma 41) and the perturbation coalgebra-morphism argument (Lemma 43) are both independently verified correct by the technical\_correctness specialist.
  \item Lemma 24's proof that every fibration in L\ensuremath{\infty} Alg(R)\_fgp is strict up to isomorphism is a technically precise application of the fgp condition and filtration completeness, and correctly underpins Proposition 23's fiber-product construction.
  \item The paper clearly situates its main result as extending the Behrend--Liao--Xu and Carchedi base CFO results to the fibered category of L\ensuremath{\infty} spaces, and explicitly distinguishes its stricter contraction notion from the Amorim--Tu relaxed approach.
\end{itemize}

\section*{Weaknesses}
\begin{itemize}[leftmargin=*]
  \item The headline theorem (Theorem 1/30) ships no machine-checkable formal proof artifact; the technical\_correctness specialist rated it partially\_supported at major severity and the reproducibility specialist flagged it at critical severity --- together triggering the math.* recommendation gate for proof-as-code.
  \item The abstract's corollary that transitive L\ensuremath{\infty} algebroids form a CFO is contingent on an unpublished companion paper (cattaneojiang26, 2026, no arXiv ID) that was unavailable for review; the dependency is not marked conditional in the abstract or theorem statement.
  \item Proposition 15's proof delegates Lemma 16 (kernel of a surjection of locally free O\_M-modules is locally free) to the same unavailable companion paper (Proposition 2.6 of cattaneojiang26), leaving a gap in the verification of the [a2] fibrancy assumption.
  \item Three bodies of directly relevant prior art are absent from the bibliography: Q-algebroids and dg Lie algebroids (Mehta, Vaintrob, Kontsevich school), Abad--Crainic representations up to homotopy, and Nuiten's Koszul duality for Lie algebroids; their omission obscures the paper's positioning within the derived geometry literature.
  \item The Serre--Swan generalization to graded ringed spaces (Lemmas 12--13) applies Morye's Lemma 2.3 and Proposition 2.5 --- stated for ungraded ringed spaces --- to the graded setting with only a footnote asserting the extension is straightforward, without verification or a graded-setting citation.
\end{itemize}

\section*{Revision Targets}
\begin{itemize}[leftmargin=*]
  \item [open] paper\_tex at `Section 1 (Theorem 1) and Section 3 (Theorem 30).`: Per the proof-as-code axiom for math.*: provide a formal verification of the CFO axioms (terminal object, pullback of fibrations, path-space factorization) and of the framing functor, e.g. in src/proofs/CFO\_LinftySp.lean (Lean 4 / mathlib4), so a machine-checkable artifact certifies Theorem 1/30. Check: Re-review should confirm `Section 1 (Theorem 1) and Section 3 (Theorem 30).` is corrected or justified.
  \item [open] paper\_tex at `Abstract; introduction.`: Make the dependency on cattaneojiang26 explicit in the abstract and theorem statement, and either include the relevant equivalence statement as an assumption or wait for the companion to be publicly available before quoting the corollary. Check: Re-review should confirm `Abstract; introduction.` is corrected or justified.
  \item [open] paper\_tex at `Section 2 (Proposition 15), proof using Lemma 17 and Lemma 16.`: Either include a self-contained proof of Lemma 16 or, since cattaneojiang26 is cited as the dependency, gate this paper on the companion's availability and explicitly state the assumption (e.g. parity of grading or characteristic) under which Lemma 16 is proved. Check: Re-review should confirm `Section 2 (Proposition 15), proof using Lemma 17 and Lemma 16.` is corrected or justified.
  \item [open] bibliography at `Work on Q-algebroids / dg Lie algebroids (e.g., Mehta, Vaintrob, or Kontsevich-school)`: Add or discuss missing prior art `Work on Q-algebroids / dg Lie algebroids (e.g., Mehta, Vaintrob, or Kontsevich-school)`. The paper's ultimate goal is a homotopy theory for L\_\ensuremath{\infty} algebroids over dg manifolds, yet it cites no prior literature on dg Lie algebroids or Q-algebroids. References to Mehta's thesis/papers or related work on infinitesimal symmetries in derived geometry would situate the motivation more precisely. Check: Re-review should confirm the related-work discussion addresses this prior art.
  \item [open] paper\_tex at `Section 2 (Lemmas 12, 13, Proposition 14).`: Briefly verify (or cite a precise reference for) the graded-ringed-space generalization of Morye's Lemma 2.3 / Proposition 2.5, since the cited paper is stated for the ungraded setting. Check: Re-review should confirm `Section 2 (Lemmas 12, 13, Proposition 14).` is corrected or justified.
\end{itemize}

\section*{Open Questions}
\begin{itemize}[leftmargin=*]
  \item Can the CFO axioms for Theorem 1/30 be formalized in Lean 4 / mathlib4, or can the authors supply a detailed proof-gap analysis identifying every non-obvious step and the precise lemma in Getzler or Rogers from which it follows?
  \item What is the expected public availability of cattaneojiang26, and would it be feasible to include Lemma 16 (kernel freeness) and a statement of the algebroid--space equivalence as a self-contained appendix so that this paper stands independently?
  \item How does the homotopy theory developed here relate to Q-algebroids and dg Lie algebroids in the sense of Mehta and Vaintrob --- are transitive L\ensuremath{\infty} algebroids over a dg manifold a special case, or does the relationship require additional structure?
  \item Does the graded-ringed-space generalization of Morye's Lemma 2.3 / Proposition 2.5 involve any sign, parity, or completeness subtleties beyond the ungraded case, or is there a reference in the literature where this exact generalization is established?
  \item Can the Dupont contraction anticommutation identities h\_n\^{}i h\_n\^{}j + h\_n\^{}j h\_n\^{}i = 0 and epsilon\_i\^{}n h\_n\^{}i = 0 used in Proposition 27 be verified explicitly within this paper rather than imported from Getzler 2009 Lemma 3.4 / Theorem 3.11 without reproof?
\end{itemize}

\section*{Per-Agent Reviews}
\subsection*{citation (\texttt{claude-sonnet-4-6})}
\begin{quote}\small\ttfamily\raggedright
\{\\
  "confidence": 0.0,\\
  "entries": [],\\
  "missing\_references": [],\\
  "summary": "Citation-use agent failed before producing a normal citation relevance review. Deterministic citation verification still runs separately; see external citation checks and verifier provenance for existence, DOI, URL, and resolver evidence. Failure: agent \textbackslash{}"citation\textbackslash{}" timed out after 180s at supervisor level"\\
\}\\
\end{quote}

\subsection*{meta\_reviewer (\texttt{claude-sonnet-4-6})}
\begin{quote}\small\ttfamily\raggedright
\{\\
  "confidence": 0.72,\\
  "questions": [\\
    "Can the CFO axioms for Theorem 1/30 be formalized in Lean 4 / mathlib4, or can the authors supply a detailed proof-gap analysis identifying every non-obvious step and the precise lemma in Getzler or Rogers from which it follows?",\\
    "What is the expected public availability of cattaneojiang26, and would it be feasible to include Lemma 16 (kernel freeness) and a statement of the algebroid--space equivalence as a self-contained appendix so that this paper stands independently?",\\
    "How does the homotopy theory developed here relate to Q-algebroids and dg Lie algebroids in the sense of Mehta and Vaintrob --- are transitive L\ensuremath{\infty} algebroids over a dg manifold a special case, or does the relationship require additional structure?",\\
    "Does the graded-ringed-space generalization of Morye's Lemma 2.3 / Proposition 2.5 involve any sign, parity, or completeness subtleties beyond the ungraded case, or is there a reference in the literature where this exact generalization is established?",\\
    "Can the Dupont contraction anticommutation identities h\_n\^{}i h\_n\^{}j + h\_n\^{}j h\_n\^{}i = 0 and epsilon\_i\^{}n h\_n\^{}i = 0 used in Proposition 27 be verified explicitly within this paper rather than imported from Getzler 2009 Lemma 3.4 / Theorem 3.11 without reproof?"\\
  ],\\
  "recommendation": "major\_revision",\\
  "revision\_targets": [\\
    \{\\
      "evidence": "The proof reduces the claim to Proposition 28 (CFO structure on L\_infty Alg(R\_M)\_fgp) via the equivalence in Proposition 18, then transports the simplicial-frame and pullback constructions back along Gamma. The argument is internally coherent and rests on Propositions 15, 23, 27, which are all proved in the paper, but the result is a load-bearing, code-amenable categorical statement that ships no executable formal proof.",\\
      "id": "weakness-1",\\
      "locator": "Section 1 (Theorem 1) and Section 3 (Theorem 30).",\\
      "required\_update": "Per the proof-as-code axiom for math.*: provide a formal verification of the CFO axioms (terminal object, pullback of fibrations, path-space factorization) and of the framing functor, e.g. in src/proofs/CFO\_LinftySp.lean (Lean 4 / mathlib4), so a machine-checkable artifact certifies Theorem 1/30.",\\
      "source\_path": null,\\
      "source\_role": "technical\_correctness",\\
      "status": "open",\\
      "target\_kind": "paper\_tex",\\
      "verification\_check": "Re-review should confirm `Section 1 (Theorem 1) and Section 3 (Theorem 30).` is corrected or justified.",\\
      "weakness\_index": 0\\
    \},\\
    \{\\
      "evidence": "This is presented as an immediate consequence of Theorem 1 together with the first main result of the companion paper cattaneojiang26, which is not yet available for verification here (cited with date=2026 and no arXiv id). The transport of the CFO structure along an equivalence is formally trivial, but the substantive content --- the equivalence L\_infty Algd\_fib \~{}= L\_infty Sp --- is in the companion.",\\
      "id": "weakness-2",\\
      "locator": "Abstract; introduction.",\\
      "required\_update": "Make the dependency on cattaneojiang26 explicit in the abstract and theorem statement, and either include the relevant equivalence statement as an assumption or wait for the companion to be publicly available before quoting the corollary.",\\
      "source\_path": null,\\
      "source\_role": "technical\_correctness",\\
      "status": "open",\\
      "target\_kind": "paper\_tex",\\
      "verification\_check": "Re-review should confirm `Abstract; introduction.` is corrected or justified.",\\
      "weakness\_index": 1\\
    \},\\
    \{\\
      "evidence": "Fineness follows from the splitting of graded manifolds (hatOmega\_M is a C\^{}infty\_M-module, hence fine). The [a2] argument tensors a surjection rho : hatOmega tensor E -> hatOmega tensor F up from its zeroth component rho\_0, splits rho\_0 via Lemma 16, then inverts id + \~{}s o (rho - \~{}rho\_0). The invertibility step uses completeness of the filtration on hatOmega\_M, which is given. The crucial input Lemma 16 (kernel of a surjection of locally free O\_M-modules is locally free) is attributed to Proposition 2.6 of the companion paper cattaneojiang26 and not reproved here.",\\
      "id": "weakness-3",\\
      "locator": "Section 2 (Proposition 15), proof using Lemma 17 and Lemma 16.",\\
      "required\_update": "Either include a self-contained proof of Lemma 16 or, since cattaneojiang26 is cited as the dependency, gate this paper on the companion's availability and explicitly state the assumption (e.g. parity of grading or characteristic) under which Lemma 16 is proved.",\\
      "source\_path": null,\\
      "source\_role": "technical\_correctness",\\
      "status": "open",\\
      "target\_kind": "paper\_tex",\\
      "verification\_check": "Re-review should confirm `Section 2 (Proposition 15), proof using Lemma 17 and Lemma 16.` is corrected or justified.",\\
      "weakness\_index": 2\\
    \},\\
    \{\\
      "evidence": "The paper's ultimate goal is a homotopy theory for L\_\ensuremath{\infty} algebroids over dg manifolds, yet it cites no prior literature on dg Lie algebroids or Q-algebroids. References to Mehta's thesis/papers or related work on infinitesimal symmetries in derived geometry would situate the motivation more precisely.",\\
      "id": "weakness-4",\\
      "locator": "Work on Q-algebroids / dg Lie algebroids (e.g., Mehta, Vaintrob, or Kontsevich-school)",\\
      "required\_update": "Add or discuss missing prior art `Work on Q-algebroids / dg Lie algebroids (e.g., Mehta, Vaintrob, or Kontsevich-school)`. The paper's ultimate goal is a homotopy theory for L\_\ensuremath{\infty} algebroids over dg manifolds, yet it cites no prior literature on dg Lie algebroids or Q-algebroids. References to Mehta's thesis/papers or related work on infinitesimal symmetries in derived geometry would situate the motivation more precisely.",\\
      "source\_path": null,\\
      "source\_role": "novelty",\\
      "status": "open",\\
      "target\_kind": "bibliography",\\
      "verification\_check": "Re-review should confirm the related-work discussion addresses this prior art.",\\
      "weakness\_index": 3\\
    \},\\
    \{\\
      "evidence": "Lemma 12 (S fully faithful on Fgp(A)) is a clean tensor-with-finitely-generated-projective-commutes-with-limits argument. Lemma 13 relies on Lemma 2.3 and Proposition 2.5 of morye2013note for the iso S(Gamma(E)) \~{}= E, which the authors note works in the graded setting because the extension is 'straightforward' (footnote 2). The Sardanashvily refinement-of-cover argument is sketched but correct.",\\
      "id": "weakness-5",\\
      "locator": "Section 2 (Lemmas 12, 13, Proposition 14).",\\
      "required\_update": "Briefly verify (or cite a precise reference for) the graded-ringed-space generalization of Morye's Lemma 2.3 / Proposition 2.5, since the cited paper is stated for the ungraded setting.",\\
      "source\_path": null,\\
      "source\_role": "technical\_correctness",\\
      "status": "open",\\
      "target\_kind": "paper\_tex",\\
      "verification\_check": "Re-review should confirm `Section 2 (Lemmas 12, 13, Proposition 14).` is corrected or justified.",\\
      "weakness\_index": 4\\
    \}\\
  ],\\
  "strengths": [\\
    "The geometric-to-algebraic reduction via the global sections functor \ensuremath{\Gamma} is clean and rigorously supported: Proposition 18's proof correctly chains Proposition 15 and Lemma 13 to establish full faithfulness and the quasi-inverse equivalence, with the finite-rank projectivity condition doing essential work throughout.",\\
    "The homotopy transfer theorem for curved L\ensuremath{\infty} algebras over filtered cdgas (Proposition 44) is developed with explicit algebraic detail --- the Berglund symmetrized tensor homotopy identity (Lemma 41) and the perturbation coalgebra-morphism argument (Lemma 43) are both independently verified correct by the technical\_correctness specialist.",\\
    "Lemma 24's proof that every fibration in L\ensuremath{\infty} Alg(R)\_fgp is strict up to isomorphism is a technically precise application of the fgp condition and filtration completeness, and correctly underpins Proposition 23's fiber-product construction.",\\
    "The paper clearly situates its main result as extending the Behrend--Liao--Xu and Carchedi base CFO results to the fibered category of L\ensuremath{\infty} spaces, and explicitly distinguishes its stricter contraction notion from the Amorim--Tu relaxed approach."\\
  ],\\
  "summary": "The paper establishes that L\ensuremath{\infty} spaces over a dg manifold form a category of fibrant objects (CFO), with the stated corollary that transitive L\ensuremath{\infty} algebroids inherit this structure via the companion paper equivalence. The technical\_correctness specialist (confidence 0.70, overall: mostly\_sound) found the supporting scaffolding --- Lemmas 24, 43; Propositions 23, 44; Lemma 17; Lemma 41 --- fully supported on detailed inspection. The headline Theorem 1/30 was assessed as partially\_supported at major severity not because the mathematical reasoning is defective but because no machine-checkable proof artifact accompanies it. The reproducibility specialist (confidence 0.87, score 0.24) independently flags the same gap at critical severity. These two specialist findings trigger the recommendation gate: the field is math.* (code-amenable), and both technical\_correctness and reproducibility flag a missing proof-as-code for the headline claim at major or critical severity, so the default recommendation is major\_revision. Two additional cross-cutting concerns compound this: (a) Theorem 1's abstract corollary and the intermediate Lemma 16 both depend on the companion paper cattaneojiang26, which carries a 2026 date and no arXiv identifier and was therefore unavailable for any reviewer; (b) the novelty specialist (confidence 0.62, score 0.45, verdict: incremental) identified three bodies of closely related prior art absent from the bibliography. The citation specialist timed out entirely before producing output; this is a pipeline failure and not evidence of a paper-level bibliography deficiency. Specialists are not in disagreement on substantive points --- technical\_correctness and reproducibility align on the proof-as-code gap, and novelty's incremental verdict is consistent with technical\_correctness finding a sound but not transformative extension of existing frameworks.",\\
  "weaknesses": [\\
    "The headline theorem (Theorem 1/30) ships no machine-checkable formal proof artifact; the technical\_correctness specialist rated it partially\_supported at major severity and the reproducibility specialist flagged it at critical severity --- together triggering the math.* recommendation gate for proof-as-code.",\\
    "The abstract's corollary that transitive L\ensuremath{\infty} algebroids form a CFO is contingent on an unpublished companion paper (cattaneojiang26, 2026, no arXiv ID) that was unavailable for review; the dependency is not marked conditional in the abstract or theorem statement.",\\
    "Proposition 15's proof delegates Lemma 16 (kernel of a surjection of locally free O\_M-modules is locally free) to the same unavailable companion paper (Proposition 2.6 of cattaneojiang26), leaving a gap in the verification of the [a2] fibrancy assumption.",\\
    "Three bodies of directly relevant prior art are absent from the bibliography: Q-algebroids and dg Lie algebroids (Mehta, Vaintrob, Kontsevich school), Abad--Crainic representations up to homotopy, and Nuiten's Koszul duality for Lie algebroids; their omission obscures the paper's positioning within the derived geometry literature.",\\
    "The Serre--Swan generalization to graded ringed spaces (Lemmas 12--13) applies Morye's Lemma 2.3 and Proposition 2.5 --- stated for ungraded ringed spaces --- to the graded setting with only a footnote asserting the extension is straightforward, without verification or a graded-setting citation."\\
  ]\\
\}\\
\end{quote}

\subsection*{novelty (\texttt{claude-sonnet-4-6})}
\begin{quote}\small\ttfamily\raggedright
\{\\
  "confidence": 0.62,\\
  "missing\_prior\_art": [\\
    \{\\
      "reason": "The paper's ultimate goal is a homotopy theory for L\_\ensuremath{\infty} algebroids over dg manifolds, yet it cites no prior literature on dg Lie algebroids or Q-algebroids. References to Mehta's thesis/papers or related work on infinitesimal symmetries in derived geometry would situate the motivation more precisely.",\\
      "title": "Work on Q-algebroids / dg Lie algebroids (e.g., Mehta, Vaintrob, or Kontsevich-school)"\\
    \},\\
    \{\\
      "reason": "Representations up to homotopy and the adjoint representation of a Lie algebroid are closely related to transitive L\_\ensuremath{\infty} algebroids; a citation would clarify how the objects studied here relate to classical Lie algebroid homotopy theory.",\\
      "title": "Abad--Crainic: Representations up to homotopy"\\
    \},\\
    \{\\
      "reason": "Nuiten's work on formal moduli problems and Koszul duality for Lie algebroids over derived bases overlaps conceptually with L\_\ensuremath{\infty} algebroids over dg manifolds; its omission may obscure the relationship between the present results and the formal moduli perspective.",\\
      "title": "Nuiten: Koszul duality for Lie algebroids"\\
    \}\\
  ],\\
  "novelty\_score": 0.45,\\
  "related\_work": [\\
    \{\\
      "citation\_key": "costello2011geometric",\\
      "delta": "Costello introduced L\_\ensuremath{\infty} spaces over smooth manifolds and their weak equivalences; this paper adapts that notion and the associated homotopy theory to the strictly more general setting of dg manifolds.",\\
      "relation": "prior\_art",\\
      "title": "A geometric construction of the Witten genus, II"\\
    \},\\
    \{\\
      "citation\_key": "getzler2025higherholonomycurvedlinftyalgebras",\\
      "delta": "This paper directly adapts Getzler's constructions---pullbacks of fibrations, simplicial frames via Dupont contraction, and the filtration-compatible treatment of curved L\_\ensuremath{\infty} algebras---from the field/pro-nilpotent setting to finitely generated projective modules over filtered cdgas, which is the algebraic core needed for the dg-manifold geometry.",\\
      "relation": "builds\_on",\\
      "title": "Higher holonomy for curved L\_\ensuremath{\infty}-algebras 1: simplicial methods"\\
    \},\\
    \{\\
      "citation\_key": "getzler2009",\\
      "delta": "Getzler's foundational Lie theory for L\_\ensuremath{\infty} algebras (polynomial de Rham forms, Dupont contraction side conditions) is the template that the current paper adapts to the curved, filtered-cdga setting; Lemma 25 and the Whitney-complex retract are taken directly from this source.",\\
      "relation": "prior\_art",\\
      "title": "Lie theory for nilpotent L\_\ensuremath{\infty}-algebras"\\
    \},\\
    \{\\
      "citation\_key": "berglund2014homological",\\
      "delta": "The symmetrized tensor homotopy (Berglund's 'h\^{}\ensuremath{\Sigma}') used in the appendix HPT is cited directly; this paper re-derives it for curved L\_\ensuremath{\infty} algebras over filtered cdgas and identifies it as the harmonization of L\_h.",\\
      "relation": "builds\_on",\\
      "title": "Homological perturbation theory for algebras over operads"\\
    \},\\
    \{\\
      "citation\_key": "rogers2020explicit",\\
      "delta": "Rogers' explicit construction of path objects and pullbacks for L\_\ensuremath{\infty} algebras over a field is adapted; the paper cites Rogers for the proof structure of Proposition 23 (pullbacks of fibrations).",\\
      "relation": "builds\_on",\\
      "title": "An explicit model for the homotopy theory of finite-type Lie n-algebras"\\
    \},\\
    \{\\
      "citation\_key": "rogers2023complete",\\
      "delta": "Rogers' homotopy theory for complete L\_\ensuremath{\infty} algebras over a field supplies the foundational CFO framework that is extended here to the filtered cdga / dg-manifold setting.",\\
      "relation": "prior\_art",\\
      "title": "Complete L\_\ensuremath{\infty}-algebras and their homotopy theory"\\
    \},\\
    \{\\
      "citation\_key": "rogers2020homotopy",\\
      "delta": "Rogers-Zhu show that L\_\ensuremath{\infty} groupoids in the Banach manifold setting form an *incomplete* CFO; the current paper obtains a full CFO for L\_\ensuremath{\infty} spaces over a dg manifold, filling this gap.",\\
      "relation": "prior\_art",\\
      "title": "On the homotopy theory for Lie\ensuremath{\infty}-groupoids, with an application to integrating L\ensuremath{\infty}-algebras"\\
    \},\\
    \{\\
      "citation\_key": "brown1973abstract",\\
      "delta": "Brown's original definition of a category of fibrant objects is the central structure this paper establishes for L\_\ensuremath{\infty} spaces over dg manifolds; no modification to the axioms is made.",\\
      "relation": "prior\_art",\\
      "title": "Abstract homotopy theory and generalized sheaf cohomology"\\
    \},\\
    \{\\
      "citation\_key": "Behrend2020thx",\\
      "delta": "Behrend-Liao-Xu proved dg manifolds form a CFO; the paper situates its main theorem as extending this base-CFO result to the fibered category of L\_\ensuremath{\infty} spaces over dg manifolds.",\\
      "relation": "prior\_art",\\
      "title": "Derived Differentiable Manifolds"\\
    \},\\
    \{\\
      "citation\_key": "carchedi2023derivedmanifoldsdifferentialgraded",\\
      "delta": "Carchedi independently proved dg manifolds form a CFO; cited alongside Behrend et al. as justification for the expectation that a total CFO structure exists.",\\
      "relation": "prior\_art",\\
      "title": "Derived Manifolds as Differential Graded Manifolds"\\
    \},\\
    \{\\
      "citation\_key": "amorim2022inverse",\\
      "delta": "Amorim-Tu study homotopy transfers and strong deformation retracts for curved L\_\ensuremath{\infty} spaces over smooth manifolds using a relaxed side condition (dropping [h,r]=0); the current paper opts for a stricter contraction notion and works over dg manifolds, giving a complementary result.",\\
      "relation": "prior\_art",\\
      "title": "The inverse function theorem for curved L-infinity spaces."\\
    \},\\
    \{\\
      "citation\_key": "cattaneojiang26",\\
      "delta": "This is the companion paper (Part I of the same program) establishing the equivalence between transitive L\_\ensuremath{\infty} algebroids and L\_\ensuremath{\infty} spaces over dg manifolds; together the two papers imply the CFO result for L\_\ensuremath{\infty} algebroids.",\\
      "relation": "orthogonal",\\
      "title": "From L\_\ensuremath{\infty} algebroids to L\_\ensuremath{\infty} spaces: Part I"\\
    \},\\
    \{\\
      "citation\_key": "dupont1976simplicial",\\
      "delta": "Dupont's contraction (and the Whitney complex) is the key analytic input for constructing simplicial frames in L\_\ensuremath{\infty} Alg(R)\_fgp; Lemma 25 records the side conditions needed for the HPT.",\\
      "relation": "builds\_on",\\
      "title": "Simplicial de Rham cohomology and characteristic classes of flat bundles"\\
    \},\\
    \{\\
      "citation\_key": "morye2013note",\\
      "delta": "Morye's Lemma 2.3 and Proposition 2.5 are invoked to prove that the global sections functor \ensuremath{\Gamma} is fully faithful for the graded ringed space (M, \ensuremath{\Omega}\^{}\_M), a technical step that underpins the reduction from geometric to algebraic homotopy theory.",\\
      "relation": "builds\_on",\\
      "title": "Note on the Serre-Swan theorem"\\
    \}\\
  ],\\
  "verdict": "incremental"\\
\}\\
\end{quote}

\subsection*{reproducibility (\texttt{gpt-5.5})}
\begin{quote}\small\ttfamily\raggedright
\{\\
  "code\_availability": "unspecified",\\
  "code\_url": null,\\
  "concerns": [\\
    \{\\
      "area": "code",\\
      "description": "No public code repository, formal proof artifact, license, or pinned release is provided for reproducing the headline theorem that L\_infinity spaces over a dg manifold form a category of fibrant objects; a formal artifact such as formalization/Theorem1\_CFO.lean would close this proof-as-code gap.",\\
      "severity": "critical"\\
    \},\\
    \{\\
      "area": "code",\\
      "description": "The homotopy transfer result for curved L\_infinity algebras over filtered cdgas is load-bearing for the paper but is only supplied as mathematical prose; a machine-checkable proof such as formalization/HomotopyTransfer.lean is not provided.",\\
      "severity": "major"\\
    \},\\
    \{\\
      "area": "code",\\
      "description": "The abstract's implication for transitive L\_infinity algebroids depends on the companion paper result, but no bundled formalization or reproducible proof artifact links that dependency to this paper; an artifact such as formalization/CompanionEquivalence.lean would be needed for end-to-end reproduction.",\\
      "severity": "major"\\
    \},\\
    \{\\
      "area": "other",\\
      "description": "No computational environment, proof-assistant version, dependency lockfile, or build instructions are specified because the paper is presented as a pure mathematical proof rather than a reproducible formal or numerical artifact.",\\
      "severity": "major"\\
    \}\\
  ],\\
  "confidence": 0.87,\\
  "data\_availability": "unspecified",\\
  "data\_url": null,\\
  "environment": \{\\
    "dependencies": [\\
      "Mathematical setting: filtered commutative differential graded algebras over R with descending, separated, exhaustive, complete filtrations",\\
      "Curved L\_infinity algebras and curved L\_infinity spaces over dg manifolds",\\
      "Homological perturbation theory and category-of-fibrant-objects axioms"\\
    ],\\
    "hardware": null,\\
    "software": null\\
  \},\\
  "reproducibility\_score": 0.24\\
\}\\
\end{quote}

\subsection*{summary (\texttt{claude-haiku-4-5})}
\begin{quote}\small\ttfamily\raggedright
\{\\
  "audience": "Researchers in differential geometry, algebraic topology, and higher algebra, particularly those working with derived geometry, dg manifolds, L\ensuremath{\infty} structures, and homotopical algebra.",\\
  "key\_contributions": [\\
    "Proves that L\ensuremath{\infty} spaces over dg manifolds form a category of fibrant objects (Theorem 1)",\\
    "Establishes that the fibrancy functor from algebroids to spaces can be interpreted as a fibrant replacement",\\
    "Develops a homotopy transfer theorem for curved L\ensuremath{\infty} algebras over filtered cdgas",\\
    "Shows that the global sections functor from L\ensuremath{\infty} spaces to curved L\ensuremath{\infty} algebras is fully faithful and establishes an equivalence with a suitable subcategory",\\
    "Provides technical machinery for comparing homotopy theories of geometric and algebraic structures via homological perturbation of curved complexes"\\
  ],\\
  "plain\_language\_summary": "This paper studies the homotopy theory of L\ensuremath{\infty} spaces---geometric structures that generalize Lie algebras to allow higher-order bracket operations---over differential graded (dg) manifolds. The main result establishes that these spaces form a category of fibrant objects, meaning they have the right structure for homotopy-theoretic computations. This extends classical homotopy transfer theorems, which describe how algebraic structures can be transferred along deformation retracts, to the curved setting where the underlying differential has non-zero square. The authors develop this theory by translating geometric questions about L\ensuremath{\infty} spaces into algebraic problems about curved L\ensuremath{\infty} algebras over filtered commutative differential graded algebras, then applying adapted versions of existing homotopy transfer techniques.",\\
  "tldr": "The paper proves that L\ensuremath{\infty} spaces over dg manifolds form a category of fibrant objects, extending homotopy transfer theorems for curved L\ensuremath{\infty} algebras to the filtered setting."\\
\}\\
\end{quote}

\subsection*{technical\_correctness (\texttt{claude-opus-4-7})}
\begin{quote}\small\ttfamily\raggedright
\{\\
  "claims": [\\
    \{\\
      "assessment": "partially\_supported",\\
      "claim": "Theorem 1 (= Theorem 30): The category L\_infty Sp(M) of curved L\_infty spaces over a dg manifold M forms a category of fibrant objects (CFO).",\\
      "evidence": "The proof reduces the claim to Proposition 28 (CFO structure on L\_infty Alg(R\_M)\_fgp) via the equivalence in Proposition 18, then transports the simplicial-frame and pullback constructions back along Gamma. The argument is internally coherent and rests on Propositions 15, 23, 27, which are all proved in the paper, but the result is a load-bearing, code-amenable categorical statement that ships no executable formal proof.",\\
      "id": "C1",\\
      "location": "Section 1 (Theorem 1) and Section 3 (Theorem 30).",\\
      "severity": "major",\\
      "suggested\_fix": "Per the proof-as-code axiom for math.*: provide a formal verification of the CFO axioms (terminal object, pullback of fibrations, path-space factorization) and of the framing functor, e.g. in src/proofs/CFO\_LinftySp.lean (Lean 4 / mathlib4), so a machine-checkable artifact certifies Theorem 1/30."\\
    \},\\
    \{\\
      "assessment": "supported",\\
      "claim": "Proposition 18: The global sections functor Gamma : L\_infty Sp(M) -> L\_infty Alg(R\_M)\_fgp is well-defined and fully faithful, and restricts to a quasi-inverse equivalence with S on L\_infty Alg(R\_M)\_vec.",\\
      "evidence": "The proof combines Proposition 15 (verifying assumptions [a1] and [a2] for (M, hatOmega\_M)) with Lemma 13 (general Serre--Swan-style equivalence under [a1], [a2]). The argument that S(g\_M) coincides with the completed base change along restriction maps and is therefore a sheaf of curved L\_infty algebras is correct given finite-rank projectivity. The Sardanashvily refinement-of-cover argument is standard.",\\
      "id": "C2",\\
      "location": "Section 2 (Proposition 18), proof immediately after.",\\
      "severity": "info",\\
      "suggested\_fix": null\\
    \},\\
    \{\\
      "assessment": "supported",\\
      "claim": "Proposition 23: For a fibration phi : g -> m and arbitrary psi : n -> m in L\_infty Alg(R)\_fgp, the fiber product g x\_m n exists, and if phi is a trivial fibration then so is the projection \~{}phi : g x\_m n -> n.",\\
      "evidence": "After reducing to the strict case via Lemma 24, the explicit construction on k oplus n (where k = ker phi\_1) with the brackets characterized by eqs. (compphi)/(pullback-morphism) is consistent. The unique-solution claim for \~{}psi was spot-checked: with p = id - s phi\_1, one has p s = 0, so p (x\_1 + s psi\_1(y\_1)) = x\_1 and p (s psi\_n) = 0, matching the stated side conditions. The D\^{}2 = 0 argument via injectivity of (CE(\~{}phi), CE(\~{}psi)) is the standard Getzler/Rogers argument cited.",\\
      "id": "C3",\\
      "location": "Section 3 (Proposition 23); proof spans pages following Lemma 24.",\\
      "severity": "info",\\
      "suggested\_fix": null\\
    \},\\
    \{\\
      "assessment": "supported",\\
      "claim": "Lemma 24: Every fibration in L\_infty Alg(R)\_fgp is strict up to isomorphism.",\\
      "evidence": "The splitting s : m[1] -> g[1] of phi\_1 exists because m is projective over R\^{}sharp (this is the substantive use of the fgp condition). The conjugating automorphism Psi has Psi\_n = s o phi\_n for n != 1, and the computed Taylor coefficients of Phi o Psi reduce to phi\_n by phi\_1 o s = id. The check that Psi is invertible uses completeness of the filtration on hatSym(g[1]), which holds.",\\
      "id": "C4",\\
      "location": "Section 3 (Lemma 24).",\\
      "severity": "info",\\
      "suggested\_fix": null\\
    \},\\
    \{\\
      "assessment": "partially\_supported",\\
      "claim": "Proposition 27: For each n, the diagram l\_n g -> g\_n = W\_n\^{}sharp tensor g ->> r\_n g = g\^{}\{oplus(n+1)\} is a factorization of the diagonal in L\_infty Alg(R)\_fgp, with the left leg a weak equivalence and the right leg a fibration.",\\
      "evidence": "The weak-equivalence leg uses acyclicity of Omega\_n and the homotopy transfer along Dupont's contraction (Lemmas 25, 26 + Proposition 44). The right-leg surjectivity reduces to surjectivity of vertex evaluation on W\_n\^{}sharp, which is standard. Commutativity of the diagram is reduced to s\_n o iota\_n = 0 and ev\_n o s\_n = 0, both consequences of h\_n\^{}i(1) = 0 and epsilon\_i\^{}n h\_n\^{}i = 0 plus anticommutation h\_n\^{}i h\_n\^{}j + h\_n\^{}j h\_n\^{}i = 0; these identities are quoted from Getzler 2009 (Lemma 3.4, Theorem 3.11) rather than reproved.",\\
      "id": "C5",\\
      "location": "Section 3 (Proposition 27).",\\
      "severity": "minor",\\
      "suggested\_fix": "Either include the explicit verification of the Dupont side conditions (h\_n\^{}i h\_n\^{}j + h\_n\^{}j h\_n\^{}i = 0 and epsilon\_i\^{}n h\_n\^{}i = 0) or, per the proof-as-code axiom, ship a checked computation (e.g. experiments/dupont/contraction\_identities.py) for the Whitney/Dupont identities used in this proof."\\
    \},\\
    \{\\
      "assessment": "supported",\\
      "claim": "Lemma 41 (Berglund): For n > 0, G restricted to Sym\^{}n\_\{R\^{}sharp\}(g[1]) equals (1/n) sum\_\{epsilon in \{0,1\}\^{}n\} C(n-1, |epsilon|)\^{}\{-1\} (ip)\^{}\{epsilon\_1\} tensor ... tensor (ip)\^{}\{epsilon\_n\}, with the convention C(n-1, n)\^{}\{-1\} = 0.",\\
      "evidence": "Spot-checked the combinatorial identity c\_\{|epsilon|\} = (1/n) [(n-|epsilon|)/C(n-1,|epsilon|) - |epsilon|/C(n-1,|epsilon|-1)]. For 0 < |epsilon| < n one has (n-|epsilon|) (n-1-|epsilon|)! = (n-|epsilon|)!, so both terms equal |epsilon|!(n-|epsilon|)!/(n-1)!, giving c\_\{|epsilon|\} = 0. The boundary values c\_0 = 1, c\_n = -1 are immediate. The result agrees with Berglund's symmetrized tensor trick.",\\
      "id": "C6",\\
      "location": "Appendix (Lemma 41) and proof.",\\
      "severity": "info",\\
      "suggested\_fix": null\\
    \},\\
    \{\\
      "assessment": "supported",\\
      "claim": "Lemma 43: After perturbation by mu = D\_g - d, the maps p\_mu and i\_mu are morphisms of filtered graded cocommutative coalgebras and D\_h = delta + p\_mu mu i is a coderivation compatible with D\_R.",\\
      "evidence": "The coalgebra-morphism identity (p\_mu tensor p\_mu) Delta = Delta p\_mu is derived using (h tensor h) Delta h = 0 and related vanishings, plus invertibility of id + mu h on the complete filtered coalgebra. The combinatorial step (h mu)\^{}k ip = (1/k!)(L\_h mu)\^{}k ip is justified using Lemma 41 and the (1/m) sum\_\{epsilon\} C(m-1,|epsilon|)\^{}\{-1\} = 1 identity, which checks: sum\_\{r=0\}\^{}\{m-1\} C(m-1,r) C(m-1,r)\^{}\{-1\} = m. The coderivation/D\_R compatibility of D\_h follows from the listed pieces being coderivations / coalgebra morphisms / R\^{}sharp-linear, as claimed.",\\
      "id": "C7",\\
      "location": "Appendix (Lemma 43) and proof.",\\
      "severity": "info",\\
      "suggested\_fix": null\\
    \},\\
    \{\\
      "assessment": "supported",\\
      "claim": "Proposition 44 (Homotopy Transfer Theorem): A contraction h of a curved L\_infty algebra g induces a curved L\_infty algebra structure on h = ker[l\_1, h] with curvature p(l\_0) and unary bracket p l\_1 i, and the perturbed maps p\_mu, i\_mu are weak equivalences with (p\_mu)\_0 = (i\_mu)\_0 = 0, (p\_mu)\_1 = p, (i\_mu)\_1 = i.",\\
      "evidence": "Follows formally from Lemma 43 plus the linear-component readoff. The weak-equivalence claim uses that (Gr p, Gr i, Gr h) is a strong deformation retract of complexes from (Gr g, Gr l\_1) to (Gr h, Gr p l\_1 i), which is immediate from the definitions and the contraction axioms. The closing computation pi(l\_0) = l\_0 and ip(l\_0) = l\_0 (using h(l\_0) = 0) confirms curvature is preserved.",\\
      "id": "C8",\\
      "location": "Appendix (Proposition 44).",\\
      "severity": "info",\\
      "suggested\_fix": null\\
    \},\\
    \{\\
      "assessment": "partially\_supported",\\
      "claim": "Proposition 15: The graded ringed space (M, hatOmega\_M) underlying a curved L\_infty space is fine ([a1]) and kernels of surjective morphisms between locally free finite-rank hatOmega\_M-modules are again locally free ([a2]).",\\
      "evidence": "Fineness follows from the splitting of graded manifolds (hatOmega\_M is a C\^{}infty\_M-module, hence fine). The [a2] argument tensors a surjection rho : hatOmega tensor E -> hatOmega tensor F up from its zeroth component rho\_0, splits rho\_0 via Lemma 16, then inverts id + \~{}s o (rho - \~{}rho\_0). The invertibility step uses completeness of the filtration on hatOmega\_M, which is given. The crucial input Lemma 16 (kernel of a surjection of locally free O\_M-modules is locally free) is attributed to Proposition 2.6 of the companion paper cattaneojiang26 and not reproved here.",\\
      "id": "C9",\\
      "location": "Section 2 (Proposition 15), proof using Lemma 17 and Lemma 16.",\\
      "severity": "minor",\\
      "suggested\_fix": "Either include a self-contained proof of Lemma 16 or, since cattaneojiang26 is cited as the dependency, gate this paper on the companion's availability and explicitly state the assumption (e.g. parity of grading or characteristic) under which Lemma 16 is proved."\\
    \},\\
    \{\\
      "assessment": "supported",\\
      "claim": "Lemma 17: Every locally free graded hatOmega\_M-module g of finite total rank is of the form hatOmega\_M tensor\_\{O\_M\} E for some locally free graded O\_M-module E of finite total rank.",\\
      "evidence": "The proof exhibits a splitting iota : E = Gr\^{}0 g -> g of the short exact sequence 0 -> F\^{}1 g -> g -> E -> 0 (using fineness and local freeness), then shows the induced \~{}iota : hatOmega tensor E -> g is locally an isomorphism. Completeness of the filtration on hatOmega\_M is invoked correctly to upgrade 'locally iso' to 'iso'.",\\
      "id": "C10",\\
      "location": "Section 2 (Lemma 17).",\\
      "severity": "info",\\
      "suggested\_fix": null\\
    \},\\
    \{\\
      "assessment": "partially\_supported",\\
      "claim": "Lemma 12 / Lemma 13 (Serre--Swan-style): For (X, O\_X) satisfying [a1] and [a2], Gamma is fully faithful, S Gamma = id on O\_X-mod, and Gamma : Vec(X) -> Fgp(A)\_vec is a quasi-inverse equivalence to S.",\\
      "evidence": "Lemma 12 (S fully faithful on Fgp(A)) is a clean tensor-with-finitely-generated-projective-commutes-with-limits argument. Lemma 13 relies on Lemma 2.3 and Proposition 2.5 of morye2013note for the iso S(Gamma(E)) \~{}= E, which the authors note works in the graded setting because the extension is 'straightforward' (footnote 2). The Sardanashvily refinement-of-cover argument is sketched but correct.",\\
      "id": "C11",\\
      "location": "Section 2 (Lemmas 12, 13, Proposition 14).",\\
      "severity": "minor",\\
      "suggested\_fix": "Briefly verify (or cite a precise reference for) the graded-ringed-space generalization of Morye's Lemma 2.3 / Proposition 2.5, since the cited paper is stated for the ungraded setting."\\
    \},\\
    \{\\
      "assessment": "partially\_supported",\\
      "claim": "Theorem 1's corollary in the abstract: Transitive L\_infty algebroids over a dg manifold also form a category of fibrant objects.",\\
      "evidence": "This is presented as an immediate consequence of Theorem 1 together with the first main result of the companion paper cattaneojiang26, which is not yet available for verification here (cited with date=2026 and no arXiv id). The transport of the CFO structure along an equivalence is formally trivial, but the substantive content --- the equivalence L\_infty Algd\_fib \~{}= L\_infty Sp --- is in the companion.",\\
      "id": "C12",\\
      "location": "Abstract; introduction.",\\
      "severity": "minor",\\
      "suggested\_fix": "Make the dependency on cattaneojiang26 explicit in the abstract and theorem statement, and either include the relevant equivalence statement as an assumption or wait for the companion to be publicly available before quoting the corollary."\\
    \}\\
  ],\\
  "confidence": 0.7,\\
  "overall\_correctness": "mostly\_sound"\\
\}\\
\end{quote}

\section*{Corrections}
% corrections rendered from corrections table; empty on first publish
No corrections have been recorded.

\section*{Bibliography}
\begin{enumerate}[leftmargin=*]
  \item Behrend2020thx: author=Behrend, Kai, author=Liao, Hsuan-Yi, author=Xu, Ping, title=Derived Differentiable Manifolds, date=2020, eprint=2006.01376, \href{https://arxiv.org/abs/2006.01376}{arXiv:2006.01376}
  \item amorim2022inverse: title=The inverse function theorem for curved L-infinity spaces., author=Amorim, Lino, author=Tu, Junwu, journal=Journal of Noncommutative Geometry, volume=16, number=4, date=2022
  \item berglund2014homological: title=Homological perturbation theory for algebras over operads, author=Berglund, Alexander, journal=Algebraic \textbackslash{}\& Geometric Topology, volume=14, number=5, pages=2511--2548, date=2014, publisher=Mathematical Sciences Publishers
  \item brown1973abstract: title=Abstract homotopy theory and generalized sheaf cohomology, author=Brown, Kenneth S., journal=Transactions of the American Mathematical Society, volume=186, pages=419--458, date=1973
  \item carchedi2023derivedmanifoldsdifferentialgraded: author=Carchedi, David, title=Derived Manifolds as Differential Graded Manifolds, date=2023, eprint=2303.11140, \href{https://arxiv.org/abs/2303.11140}{arXiv:2303.11140}
  \item cattaneojiang26: author = Cattaneo, Alberto S., author = Jiang, Shuhan, title = From \$L\_\textbackslash{}infty\$ algebroids to \$L\_\textbackslash{}infty\$ spaces: Part I, date = 2026,
  \item costello2011geometric: author=Costello, Kevin, title=A geometric construction of the Witten genus, II, date=2011, eprint=1112.0816, \href{https://arxiv.org/abs/1112.0816}{arXiv:1112.0816}
  \item dupont1976simplicial: title=Simplicial de Rham cohomology and characteristic classes of flat bundles, author=Dupont, Johan L, journal=Topology, volume=15, number=3, pages=233--245, date=1976, publisher=Pergamon
  \item getzler2009: author=Getzler, Ezra, title=Lie theory for nilpotent \$L\_\textbackslash{}infty\$-algebras, journal=Annals of Mathematics, volume=170, number=1, date=2009, pages=271--301,
  \item getzler2025higherholonomycurvedlinftyalgebras: author=Getzler, Ezra, title=Higher holonomy for curved \$L\_\textbackslash{}infty\$-algebras 1: simplicial methods, date=2025, eprint=2408.11157, \href{https://arxiv.org/abs/2408.11157}{arXiv:2408.11157}
  \item hovey2007model: title=Model Categories, author=Hovey, Mark, number=63, date=2007, publisher=American Mathematical Soc.,
  \item morye2013note: title=Note on the Serre-Swan theorem, author=Morye, Archana S, journal=Mathematische Nachrichten, volume=286, number=2-3, pages=272--278, date=2013, publisher=Wiley Online Library
  \item rogers2020explicit: title=An explicit model for the homotopy theory of finite-type Lie n--algebras, author=Rogers, Christopher L, journal=Algebraic \textbackslash{}\& Geometric Topology, volume=20, number=3, pages=1371--1429, date=2020, publisher=Mathematical Sciences Publishers
  \item rogers2020homotopy: title=On the homotopy theory for Lie\ensuremath{\infty}--groupoids, with an application to integrating L\ensuremath{\infty}--algebras, author=Rogers, Christopher L, author=Zhu, Chenchang, journal=Algebraic \textbackslash{}\& Geometric Topology, volume=20, number=3, pages=1127--1219, date=2020, publisher=Mathematical Sciences Publishers,
  \item rogers2023complete: title=Complete \$L\_\textbackslash{}infty\$-algebras and their homotopy theory, author=Rogers, Christopher L, journal=Journal of Pure and Applied Algebra, volume=227, number=10, pages=107403, date=2023, publisher=Elsevier,
  \item sardanashvily2001remark: title=Remark on the Serre-Swan theorem for non-compact manifolds, author=Sardanashvily, G, eprint=math-ph/0102016, date=2001 \href{https://arxiv.org/abs/math-ph/0102016}{arXiv:math-ph/0102016}
\end{enumerate}

% \bibliographystyle{plain}
% \bibliography{references}
\end{document}
